\documentclass[border=10pt]{standalone}
\usepackage{fontspec}
\setmainfont{Noto Serif}
\newfontfamily\bengalifont{Noto Serif Bengali}[Script=Bengali]
\usepackage[Latin,Bengali]{ucharclasses}
\setTransitionsFor{Bengali}{\begingroup\bengalifont}{\endgroup}
\usepackage{newunicodechar}
\newunicodechar{।}{{\bengalifont\char"0964}}
\newunicodechar{॥}{{\bengalifont\char"0965}}

\usepackage[table]{xcolor}
\usepackage{booktabs,tabularx,array}
\definecolor{bapilrule}{HTML}{1F2937}
\definecolor{bapilstripe}{HTML}{F3F4F6}
\renewcommand{\arraystretch}{1.28}
\setlength{\tabcolsep}{5pt}

\begin{document}
\begin{minipage}{17cm}
{\large\bfseries সারণি ৩৩।}\ {\large এআই ব্যবস্থার নির্দিষ্ট নিরাপত্তা হুমকি ও পরিকল্পিত নিয়ন্ত্রণ}\par\vspace{7pt}
\small
\rowcolors{2}{bapilstripe}{white}
\begin{tabularx}{\linewidth}{@{}>{\hsize=0.732\hsize\raggedright\arraybackslash}X>{\hsize=0.805\hsize\raggedright\arraybackslash}X>{\hsize=1.463\hsize\raggedright\arraybackslash}X@{}}
\arrayrulecolor{bapilrule}\toprule
\rowcolor{white}\textbf{হুমকি (OWASP LLM Top 10, ২০২৫)} & \textbf{BAPIL প্রেক্ষাপটে উদাহরণ} & \textbf{পরিকল্পিত নিয়ন্ত্রণ} \\
\midrule
LLM01 প্রম্পট ইনজেকশন & গৃহীত PDF-এর ভেতরে লুকানো নির্দেশনা & উদ্ধারকৃত টেক্সটকে শুধু ডেটা হিসেবে গণ্য করা; কোনো টুল চালানো নয়; শুধু উদ্ধৃতিভিত্তিক উত্তর; প্রতিটি G2-তে রেড-টিম পরীক্ষা \\
LLM02 সংবেদনশীল তথ্য প্রকাশ & উত্তরে ব্যক্তিগত তথ্য চলে আসা & প্রথমে P0 কর্পাস; গ্রহণের সময় ব্যক্তিগত তথ্য স্ক্যান; সংগ্রহভিত্তিক প্রবেশ নিয়ন্ত্রণ; লগ থেকে তথ্য অপসারণ \\
LLM03 সরবরাহ শৃঙ্খল & ক্ষতিগ্রস্ত মডেল-ওজন বা প্যাকেজ & নির্দিষ্ট সংস্করণের নির্ভরতা; সফটওয়্যার উপাদান তালিকা (SBOM); মডেলের চেকসাম ও উৎস-তথ্য; স্বাক্ষরিত রিলিজ \\
LLM04 ডেটা ও মডেল বিষক্রিয়া & কর্পাসে কারসাজিকৃত নথি প্রবেশ & অনুমোদিত উৎসের তালিকা; উৎস-হ্যাশ; পরিবর্তন শনাক্তকরণ; ফাইন-টিউনিং ডেটার মানব-বাছাই \\
LLM05 অনুপযুক্ত আউটপুট ব্যবস্থাপনা & মডেলের আউটপুট সক্রিয় কনটেন্ট হিসেবে প্রদর্শিত & সকল আউটপুট এস্কেপ করা; কাঠামোবদ্ধ আউটপুট যাচাই; সরাসরি চালনা নয় \\
LLM06 অতিরিক্ত কর্তৃত্ব & সরঞ্জামকে ব্যবস্থায় পদক্ষেপ নেওয়ার ক্ষমতা দেওয়া & শুধু-পঠন নকশা; কোনো স্বয়ংক্রিয় পদক্ষেপ নয়; যেকোনো পরিবর্তনে মানব-অনুমোদন \\
LLM07 সিস্টেম প্রম্পট ফাঁস & প্রম্পটে গোপন তথ্য রাখা & প্রম্পটে কোনো গোপন তথ্য নয়; প্রম্পট সংস্করণকৃত ও সম্ভব হলে প্রকাশিত \\
LLM08 ভেক্টর ও এমবেডিং দুর্বলতা & সূচিতে এক সংগ্রহ থেকে অন্য সংগ্রহে তথ্য ফাঁস & অনুসন্ধানের সময় সংগ্রহভিত্তিক প্রবেশ নিয়ন্ত্রণ; ব্যবহারকারী-গোষ্ঠী পৃথকীকরণ \\
LLM09 ভুল তথ্য & আত্মবিশ্বাসী কিন্তু ভুল উত্তর & বাধ্যতামূলক উদ্ধৃতি; বিরত থাকা; আস্থার মাত্রা প্রদর্শন; ব্যবহারকারী প্রশিক্ষণ \\
LLM10 সীমাহীন সম্পদ ব্যবহার & অনুসন্ধানের বন্যায় কম্পিউটিং ব্যয় বৃদ্ধি & কোটা, হার-সীমা, সময়সীমা, ব্যয় পর্যবেক্ষণ \\
\bottomrule
\end{tabularx}
\par\vspace{3pt}{\scriptsize\color{gray} উৎস: BAPIL কারিগরি ব্লুপ্রিন্ট, সংস্করণ ১.০ (২০২৬)। CC BY 4.0।}
\end{minipage}
\end{document}
